% ============================================================
% Chapter 8: Sustainability and Ethical Implications
%   Section: Equity and Bias Analysis
%     Subsection: Critical Analysis of the Impact
% Included from chapters/08-sustainability-ethics/03-equity-bias-analysis/equity-bias-analysis.tex
% ============================================================

\subsection{Critical Analysis of the Impact}

Not every bias above matters equally. Three can change what a reader
is shown or the verdict a claim receives, and remain.

\begin{itemize}
    \item \textbf{The screen and the ranking decide what is offered.} A
    story the screen leaves out never reaches an editor unless they open
    the list of what was left out; a story ranked below the first twenty
    is rarely read. Both rest on descriptions and labels written by the
    author or an AI, so the publication's view of ``constructive'' is
    applied at scale before a person sees anything. Showing what was
    left out and why is a mitigation, not a cure.
    \item \textbf{The evidence decides the verdict.} The search engines'
    coverage and the catalogue's ratings favour well-known Western
    outlets, so a claim about a project reported mainly by local or
    non-English media is more likely to come back \texttt{UNVERIFIED},
    which reads as less trustworthy even though it is a statement about
    the evidence, not the claim.
    \item \textbf{Spanish is served less well.} Fewer sources, a less
    precise claim selector and a worse pilot result in Spanish (22.5\%
    against 41.7\% on the pilot, with different label sets) mean that
    Spanish-language news is checked less reliably.
\end{itemize}

The biases in the evaluation data, by contrast, affect what this thesis
can claim rather than what a reader sees; they are reported as threats
to validity (Section~\ref{sec:threats}). The most important of them ---
labels proposed by a model --- is also the one the project can remove
before the end, by the annotator's review.
